% Generated draft: requires mathematical and source review.
\documentclass[11pt]{article}
\usepackage[margin=25mm]{geometry}
\usepackage{fontspec}
\setmainfont{Latin Modern Roman}
\usepackage{amsmath,amssymb,mathtools,unicode-math}
\setmathfont{Latin Modern Math}
\usepackage{enumitem,needspace,xurl,hyperref}
\setlength{\parindent}{0pt}
\setlength{\parskip}{5pt}
\setlength{\emergencystretch}{3em}
\newcommand{\N}{\mathbb{N}}
\newcommand{\Z}{\mathbb{Z}}
\newcommand{\Q}{\mathbb{Q}}
\newcommand{\R}{\mathbb{R}}
\newcommand{\C}{\mathbb{C}}
\newcommand{\F}{\mathbb{F}}
\newcommand{\A}{\mathbb{A}}
\newcommand{\PP}{\mathbb{P}}
\newcommand{\E}{\mathbb{E}}
\newcommand{\eps}{\varepsilon}
\newcommand{\dd}{\,\mathrm{d}}
\begin{document}
\section*{UnsolvedMath ID: 30004990}
\textbf{Source problem code:} catalogue entry
\subsection{Definitions and mathematical statement}
The supplied record conjectures dynamical comparison for every amenable minimal action of a nonamenable group on a compact metric space. It does not specify the topology or countability of the group, or define amenability and dynamical comparison. The following is a precise conditional formulation using discrete groups, topological amenability, and the open-set comparison convention stated below. These interpretations require source review.

Let \(G\) be a discrete group and let \(X\) be a nonempty compact metric space. An action \(G\curvearrowright X\) is an action by homeomorphisms, written \((g,x)\mapsto gx\). It is minimal if \(\overline{Gx}=X\) for every \(x\in X\). The group is nonamenable if there is no left-translation-invariant mean on \(\ell^\infty(G;\R)\).

Write
\[
\operatorname{Prob}(G)=\left\{p\in\ell^1(G;\R):p(h)\geq 0\text{ for all }h\in G,\quad \sum_{h\in G}p(h)=1\right\},
\qquad (g\cdot p)(h)=p(g^{-1}h).
\]
The action is topologically amenable if there is a net of norm-continuous maps \(m_i:X\to\operatorname{Prob}(G)\) such that, for every \(g\in G\),
\[
\lim_i\sup_{x\in X}\left\|m_i(gx)-g\cdot m_i(x)\right\|_1=0.
\]
Let \(M_G(X)\) be the set of \(G\)-invariant Borel probability measures on \(X\); thus \(\mu(gA)=\mu(A)\) for every Borel set \(A\subseteq X\) and every \(g\in G\).

For nonempty open sets \(U,V\subseteq X\), write \(U\prec V\) if, for every closed set \(F\subseteq U\), there exist finitely many open sets \(W_1,\ldots,W_n\subseteq X\) and elements \(g_1,\ldots,g_n\in G\) such that
\[
F\subseteq\bigcup_{j=1}^{n}W_j,\qquad g_jW_j\subseteq V\quad(1\leq j\leq n),\qquad g_jW_j\cap g_kW_k=\varnothing\quad(j\neq k).
\]
The empty family is allowed when \(F=\varnothing\). Under the convention used here, the action has dynamical comparison if, for every pair of nonempty open sets \(U,V\subseteq X\),
\[
\left[\mu(U)<\mu(V)\text{ for every }\mu\in M_G(X)\right]\ \Longrightarrow\ U\prec V.
\]
If \(M_G(X)=\varnothing\), the bracketed condition is vacuous, so this definition requires \(U\prec V\) for every pair of nonempty open sets. No existence of an invariant probability measure is assumed.

\textbf{Conjecture, conditional on these conventions.} For every nonamenable discrete group \(G\), every nonempty compact metric space \(X\), and every topologically amenable minimal action \(G\curvearrowright X\), the action has dynamical comparison.

The supplied text imposes no freeness, finite-dimensionality, or countability hypothesis. None is added here. Whether a countable discrete group or another additional hypothesis is intended remains unresolved. The original wording places ``amenable'' in parentheses, whereas the supplied reconstructed statement explicitly requires amenability; the formulation above follows the reconstructed statement.
\subsection{Short English statement}
The conjecture says that every topologically amenable minimal action of a nonamenable group on a compact metric space has dynamical comparison: whenever all invariant probability measures assign an open set \(U\) less mass than a nonempty open set \(V\), each closed subset of \(U\) can be covered by finitely many open pieces whose group translates lie disjointly inside \(V\). The discrete-group setting and the precise comparison convention used here require confirmation from the source. The supplied status is partially solved, not an independently verified conclusion.
\subsection{Sources}
\begin{itemize}
\item Supplied local UnsolvedMath registry record, permanent ID \(30004990\), exact source problem code \texttt{catalogue entry}, title: \emph{Dynamical Comparison for Amenable Minimal Nonamenable-Group Actions}. The record belongs to the Oberwolfach Reports Open Problems collection and gives proposed year \(2022\).
\item Supplied report citation: \emph{Groups and Dynamics: Topology, Measure, and Borel Structure} (2023), Oberwolfach Reports, pp.\ 107--169, DOI \texttt{10.4171/owr/2022/3}. Supplied link: \url{https://doi.org/10.4171/owr/2022/3}. The citation's year \(2023\), the DOI's \(2022\) component, and the proposed year are retained without reconciliation.
\item Supplied attribution: the original excerpt associates the conjecture with Shirly Geffen's contribution \emph{Purely infinite crossed products by non-amenable groups}, described as joint work with Eusebio Gardella, Julian Kranz, and Petr Naryshkin. The registry's separate proposer field is empty.
\item Gardella--Geffen--Kranz--Naryshkin, \emph{Classifiability of crossed products by nonamenable groups}: \url{https://arxiv.org/abs/2201.03409}. Supplied journal-version link: \url{https://doi.org/10.1515/crelle-2023-0012}.
\item Supplied literature assessment, checked on \texttt{2026-08-22}: the record labels the problem partially solved and reports dynamical comparison for amenable minimal actions of a large class of nonamenable groups, including nonamenable hyperbolic groups and several products, but not for all such groups. This is supplied dated curation derived from the OWR record, not an independently verified literature assessment; the status and the scope of the cited results require independent checking.
\item The registry labels the statement \texttt{corrected\_verified} and says it was reconstructed and checked against the official Oberwolfach report. These are supplied provenance claims; no source visit or independent verification is asserted in this draft.
\end{itemize}
\end{document}
